\section{Prospective choice and economically interpretable comparisons}\label{sec:prospective86}
Exact ranking on a completed catalogue and preference before a new workload are different claims. We connect them by a population selection bound and a separately frozen evaluation, retaining the original investment economy and every comparator.

\begin{proposition}[Population regret for a fixed service menu]\label{prop:selection86}
Condition on construction of $m$ complete services, each with original-law cost in $[0,H]$. For $n$ independent new workloads, obtain valid pathwise cost enclosures $[l_{ji},u_{ji}]$. Write $\bar l_j,\bar u_j$ for endpoint means and $q_n=H\sqrt{\log(2m/\delta)/(2n)}$. Select the service minimizing $\lambda\bar u_j+p\cdot w_j+f_j$, with a fixed tie rule. Here $\lambda\ge0$, $p\ge0$, and the nonrandom resource and fee inputs are declared separately. With probability at least $1-\delta$, simultaneously for every such price-fee vector, its expected objective regret is at most
\[
 \lambda\{2q_n+\max_j(\bar u_j-\bar l_j)\}.
\]
A separate evaluation stream with simultaneous pairwise intervals $[l_{jk},u_{jk}]$ gives cost regret at most $\max\{0,\max_k u_{\widehat j,k}\}$ for any calibration-selected service $\widehat j$. This comparison is with the best expected-cost member of the declared menu, not the original Bellman optimum; the original-optimum account remains the separate policy sandwich.
\end{proposition}

\subsection{The frozen independent evaluation}
The prospective protocol preceding the retained evaluation fixes the nine complete two-date investment controllers: the zero controller and both ReLU and Bernstein readouts, under both legacy and calibrated qualification, at acquisition sizes eight and sixteen. Coefficients, complete action tables, constraints and original continuous innovation laws are unchanged. There are 8,192 calibration bins and 131,072 distinct evaluation bins, with forty bits per coordinate. Selection uses the smallest calibration conditioned upper-endpoint mean, with lexical ties, and is fixed before evaluation.

Both independently written whole-bin evaluators run on every policy and bin; rational midpoint fixtures check their hulls. The two estimators are cash cost and analytically conditioned cash cost. Each stream has 36 unordered pairs for each estimator. The empirical-Bernstein error account used in the original study assigns $4\exp(-11)$ to each interval, yielding $288\exp(-11)<0.005$ per stream and $576\exp(-11)<0.01$ jointly. No fitting-success event is inferred. The 139,264 paired workload rows count once, rather than once per policy or per comparison. Reconstructing the frozen streams during the current build adds no independent observations.

\begin{table}[htbp]
\centering
\caption{Independent evaluation: selected fine calibrated Bernstein minus comparator}\label{tab:selection86}
\begin{tabular}{lrr}\toprule
Comparator & Lower endpoint & Upper endpoint\\\midrule
Zero & $-0.177441$ & $-0.170788$\\
Legacy ReLU, size 8 & $-0.010137$ & $-0.005786$\\
Legacy Bernstein, size 8 & $-0.087292$ & $-0.081642$\\
Legacy ReLU, size 16 & $-0.009460$ & $-0.005046$\\
Legacy Bernstein, size 16 & $-0.086152$ & $-0.080530$\\
Calibrated ReLU, size 8 & $-0.009056$ & $-0.004812$\\
Calibrated Bernstein, size 8 & $-0.004823$ & $-0.000716$\\
Calibrated ReLU, size 16 & $-0.003559$ & $0.000573$\\\bottomrule
\end{tabular}
\par\smallskip
{\footnotesize Original-law conditioned cost. The table retains the prespecified R85 evaluation checkpoint; the current reproduction regenerates it from exact integer moments. Endpoints are rounded outwards to six decimal places. Selection uses 8,192 calibration rows, not the 131,072 evaluation rows. The simultaneous evaluation family contains 72 intervals across cash and conditioned estimators. The margin is $1/256$ in normalized cost units.}
\end{table}


Calibration selects the fine calibrated Bernstein controller. In the retained evaluation, its upper expected-cost regret is less than $0.000573$, below the prespecified $1/256$ noninferiority margin. Seven of its eight comparisons have negative upper endpoints. Its comparison with fine calibrated ReLU crosses zero and supports equivalence at that margin, not strict basis superiority. The resolution comparison within calibrated Bernstein is signed. The crossed basis-by-qualification effect is a separate factorial estimand, not a change in the underlying economy. Full endpoint moments, all 144 stream-estimator contrasts, sampling and enclosure components, and all comparisons are retained with the reproduction record.

The margin is a declared acceptable normalized loss. It is not an estimated monetary adoption fee. This evaluation does not replace the historical recursive-query workload, and does not assert that the nine acquired controllers are a prospectively validated neural query proposal on every new regime or horizon. The population construction in Section~\ref{sec:population86} supplies an explicit broader guarantee under its own joint-domain hypotheses; the two empirical objects are not conflated.

\subsection{Uniform price-fee sensitivity}\label{sec:prices86}
Let $\mathcal X$ be a finite set of feasible service assignments, including the original capacity constraints. Let $\mathcal D$ be a nonempty compact polytope containing their joint normalized cost vector. For an assignment $x$, let $w_x$ be its complete resource vector, $\lambda\ge0$ the value of normalized policy loss, $p\ge0$ the resource prices, and $f$ the assignment fee vector. All coordinates refer to explicit units. Put $h_{\mathcal D}(v)=\sup_{c\in\mathcal D}c\cdot v$.

\begin{proposition}[Uniform purchase region]\label{prop:prices86}
Assignment $x$ is optimal for every cost vector in $\mathcal D$ exactly on the region
\begin{equation}\label{eq:price-region86}
 \mathcal R_x=\bigcap_{y\in\mathcal X}
 \{(\lambda,p,f):\lambda h_{\mathcal D}(x-y)
              +p\cdot(w_x-w_y)+f\cdot(x-y)\le0\}.
\end{equation}
This region is polyhedral. Rational support-function upper certificates give an independently checkable sufficient region. Uniformity over a compact polytope of price-fee inputs can be checked at its vertices. The supporting equalities give the breakpoints at which a guaranteed choice can change. A point outside the sufficient region is not automatically evidence that the candidate is inferior.
\end{proposition}

For a single selected service and comparator $k$, a simultaneous cost upper bound $u_{\widehat j,k}$ gives the transparent sufficient condition
\[
 f_{\widehat j}-f_k+p\cdot(w_{\widehat j}-w_k)
 \le-\lambda u_{\widehat j,k}.
\]
The numerical appendix prints every selected-service breakpoint, including negative ones. In particular, the fine-versus-coarse calibrated Bernstein policy comparison supports a positive extra-fee allowance even at zero computation price, whereas a comparison whose upper cost endpoint is positive cannot justify a positive fee solely from policy quality. A resource-driven choice is distinguished from a policy-cost-driven one. Application-specific prices, energy costs and monetary scales require independent calibration; none is manufactured from the normalized tolerance or the favorable timing repetition.

\paragraph{Evidence identities.}
The inherited experiments have distinct observation units. The retained R67 description records 212 service returns and 48 neural fitting warnings; the later training, validation, reuse and vector catalogue records 188 complete processes. The attained-policy factorial study instead has nine controllers and its own independent path rows. The current complete-menu calculation has 2,548 qualification certificates and no new fits. Counts from different catalogues, including the 336-service description in the advisory report, are not added together or assigned to this evaluation without a source-matched manifest. The finer-bin experiment does not erase the older adverse timing and horizon results. Every inherited active section is retained, and the reply identifies the relevant experiment for each contested assertion.
